% %%%%%%%%%%%%%%%%%%%%%%%%%%%%%%%%%%%%%%%%%%%%%%%%%%%%%%%%%%%%%%%%%%%%%%%%%%%%%%%%%
%
%	cap1.tex 
%  -----------
%
%	Neste arquivo você escreve o capítulo. Deve sempre iniciar com o ambiente
%	\chapter{Nome do Capitulo}
%
% %%%%%%%%%%%%%%%%%%%%%%%%%%%%%%%%%%%%%%%%%%%%%%%%%%%%%%%%%%%%%%%%%%%%%%%%%%%%%%%%%
%
\chapter{Introdução}
% ----------------------------------------------------------


A execução de aberturas em elementos estruturais é uma prática comum na construção civil. O cenário econômico atual demanda a execução de empreendimentos em prazos curtos, além do avanço tecnológico que torna os projetos cada vez mais complexos, sendo necessária a atuação de múltiplas equipes. Diante desse contexto, é inevitável que ocorram falhas e incompatibilidades. As aberturas são feitas nas estruturas, tornando possível a passagem de tubulações e a execução das instalações complementares. Por outro lado, também é comum que os furos sejam considerados no projeto, visando otimizar o espaço disponível na edificação.

Tendo em vista a importância do tema e a alta demanda para engenheiros civis quanto ao dimensionamento de estruturas nessas condições, o trabalho consiste no desenvolvimento de um software para o dimensionamento de vigas em concreto armado com aberturas na direçãode sua largura. Para isso, serão considerados os critérios e orientações da NBR 6118 \cite{ABNT_NBR6118_2026}, norma brasileira que trata de projeto de estruturas de concreto e da forma como as aberturas devem ser abordadas nesses projetos.

As formas mais clássicas de analisar esses casos variam desde métodos empíricos como o de \citeonline{Leonhardt1978}, que considera restrições quanto às dimensões relativas da abertura na viga, até abordagens capazes de representar situações mais gerais como os modelos de bielas e tirantes ou método dos elementos finitos. O algoritmo será desenvolvidopelo método da rigidez, também conhecido como método matricial, o qual serviu como base para o desenvolvimento do método dos elementos finitos. Este, por sua vez, é reconhecido por ser uma das melhores ferramentas para a resolução de um amplo conjunto de problemas de engenharia estrutural. Segundo \citeonline{Mendonca2019}, a formulação matricial permite compreender, de forma simplificada, conceitos do método dos elementos finitos como graus de liberdade, matriz de rigidez, processo de sobreposição, sistema de equações algébricas de equilíbrio, imposição de condições de contorno, entre outros. O trabalho surgiu não só como uma tentativa de automatizar esse tipo de projeto, mas também como uma forma de explorar aplicações para o método da rigidez.